\section{Additional Security and Accountability Measurements}
\label{app:e-security-accountability}

\noindent\textbf{Current attack construction and search.}
The current suite uses the same underlying protocol fixtures for
\textsc{SignedLog}, \textsc{SignedLog+}, and \textsc{TokenProv}.
Each episode starts from a fresh request fixture.
The attacker selects a call and one of six mutation families, repairs the
affected outer closure and settlement commitments, and signs updated records
using corrupted roles' keys.
For cross-call splicing, it replaces selected input/output and handoff hashes
with another call's fields and re-signs the provider receipt.
For deployment substitution, it changes the deployment identifier and re-signs
the receipt.
Original authorization and retained handoff records are held fixed.

An independent oracle compares candidate evidence with the original call set,
receipt fields, and integer bill; it does not use the auditor's failure label
to decide whether an attack was constructed successfully.
All 3,300 retained attack queries pass both signature-validity and
oracle-violation checks.
Queries after a rejection may switch attack family based on the exposed
first-failed-relation feedback.
Each scheme has 100 episodes across all six families, with a maximum of 16
queries per episode and early stopping after acceptance.
An episode succeeds when at least one violating candidate is accepted.
One unmodified fixture per episode is also checked, giving 100 benign controls
per scheme.

Table~\ref{tab:e-attack-families} gives accepted queries divided by actual
queries within each family.
The schemes share starting fixtures, but feedback and early stopping produce
different search trajectories and family counts.
\textsc{SignedLog+} performs an independent disclosed-record audit of signatures,
context and call links, receipt coverage, and integer accounting.
Its zero-success result is retained alongside \textsc{TokenProv}'s result.
These are CPU protocol microbenchmarks and do not evaluate attacks on live
model execution or establish a distribution-free attack success bound.

\begin{table*}[!t]
\centering
\caption{Current valid-signature adaptive attacks by family (accepted attack queries / actual queries). Family query counts differ because search adapts to feedback and stops at the first success.}
\label{tab:e-attack-families}
\footnotesize
\renewcommand{\arraystretch}{1.10}
\begin{tabular*}{\textwidth}{@{\hspace{4pt}}l@{\extracolsep{\fill}}cccc@{\hspace{4pt}}}
\toprule
\textbf{Attack family} & \textbf{SignedLog} & \textbf{SignedLog+} & \textbf{TokenProv} & \shortstack{\textbf{TokenProv first}\\\textbf{failed relation}} \\
\midrule
Aggregate bill manipulation & 18/18 & 0/264 & 0/232 & $R_{\mathrm{set}}$ \\
Receipt duplication & 13/13 & 0/267 & 0/223 & $R_{\mathrm{prov}}$ \\
Receipt omission & 22/22 & 0/278 & 0/256 & $R_{\mathrm{set}}$ \\
Observable accepted-call omission & 12/12 & 0/265 & 0/283 & $R_{\mathrm{hand}}$ \\
Cross-call evidence splicing & 14/14 & 0/251 & 0/309 & $R_{\mathrm{prov}}$ \\
Unauthorized deployment substitution & 21/21 & 0/275 & 0/297 & $R_{\mathrm{prov}}$ \\
\bottomrule
\end{tabular*}
\end{table*}


\noindent\textbf{Historical injection coverage.}
Table~\ref{tab:e-historical-attack-coverage} reports fourteen predefined
injection classes, measured separately from the current valid-signature suite.
Each attack--scheme pair has five seed fixtures and 100 repetitions per seed,
giving 500 trials per pair and 7,000 per scheme.
\textsc{TokenProv} rejects all fourteen classes; the individual-record
\textsc{SignedLog} fixtures reject nine.
The historical suite includes mutations that leave original signatures stale.
Its \textsc{SignedLog} cases use individual signed-record fixtures, while
\textsc{TokenProv} audits complete protocol fixtures.
These counts characterize coverage of the specified templates and are not
pooled with the current common-fixture adaptive comparison.

\begin{table*}[!t]
\centering
\caption{Historical protocol attack microbenchmark (detected / 500 trials for each attack and scheme). These fixtures are separate from the current valid-signature adaptive suite.}
\label{tab:e-historical-attack-coverage}
\footnotesize
\renewcommand{\arraystretch}{1.10}
\begin{tabular*}{\textwidth}{@{\hspace{4pt}}l@{\extracolsep{\fill}}lccc@{\hspace{4pt}}}
\toprule
\textbf{Attack} & \textbf{Expected relation} & \textbf{Vanilla} & \textbf{SignedLog} & \textbf{TokenProv} \\
\midrule
A1: Unauthorized call & $R_{\mathrm{auth}}$ & 0/500 & 500/500 & 500/500 \\
A2: Call insertion & $R_{\mathrm{life}}$ & 0/500 & 0/500 & 500/500 \\
A3: Observable-call omission & $R_{\mathrm{hand}}$ & 0/500 & 0/500 & 500/500 \\
A4: Call replay & $R_{\mathrm{prov}}$ & 0/500 & 0/500 & 500/500 \\
A5: Call relabeling & $R_{\mathrm{prov}}$ & 0/500 & 500/500 & 500/500 \\
A6: $H\!\rightarrow\!B$ boundary tampering & $R_{\mathrm{hand}}$ & 0/500 & 500/500 & 500/500 \\
A7: $B\!\rightarrow\!P$ boundary tampering & $R_{\mathrm{hand}}$ & 0/500 & 500/500 & 500/500 \\
A8: Model/deployment substitution & $R_{\mathrm{prov}}$ & 0/500 & 500/500 & 500/500 \\
A9: Tokenizer/input-semantics substitution & $R_{\mathrm{prov}}$ & 0/500 & 500/500 & 500/500 \\
A10: Cache or usage inflation & $R_{\mathrm{prov}}$ & 0/500 & 500/500 & 500/500 \\
A11: Wrong price version & $R_{\mathrm{prov}}$ & 0/500 & 500/500 & 500/500 \\
A12: Cross-call evidence splicing & $R_{\mathrm{prov}}$ & 0/500 & 0/500 & 500/500 \\
A13: Receipt omission or duplication & $R_{\mathrm{set}}$ & 0/500 & 0/500 & 500/500 \\
A14: Aggregate bill manipulation & $R_{\mathrm{set}}$ & 0/500 & 500/500 & 500/500 \\
\bottomrule
\end{tabular*}
\end{table*}


\noindent\textbf{Cross-request transplant procedure.}
The historical binding experiment constructs separate source and target
requests and inserts the source's first call/receipt/artifact tuple into the
target's first provider-call slot.
Both call labels are \texttt{c-0}, but the request identifiers differ.
The source and target also use independently generated keys.
Source-local verification checks the receipt under the original source key
and registration; complete auditing checks it against the target authorization
and transcript.
All full-audit trials fail at the request-ID comparison in $R_{\mathrm{prov}}$,
before subsequent provider checks.
Thus, these measurements test request-context binding, rather than independently
isolating every component of $\mathcal{R}_{\mathsf{link}}$.
Each of the three source--target configurations has 500 repetitions per variant.
The archived \texttt{w/o-R\_link} path executes the same receipt-only branch as
the source-local checker and is treated as a local-check surrogate.

\noindent\textbf{Localization procedure.}
The five stage-specific mutations are authorization-context modification,
lifecycle-root modification, handoff-boundary tampering, usage modification,
and aggregate-bill modification.
The auditor evaluates $R_{\mathrm{auth}}$, $R_{\mathrm{life}}$,
$R_{\mathrm{hand}}$, $R_{\mathrm{prov}}$, and $R_{\mathrm{set}}$ in order,
returning the first failure.
Relation agreement compares this output with the predefined injection stage.
Returned intervals are fixed by the failed relation:
$U$--$H$, $H$--$H$, $H$--$P$, $B$--$P$, and $H$--$B$, respectively.
Widths count edges along $U\!\rightarrow\!H\!\rightarrow\!B\!\rightarrow\!P$;
zero width denotes one role.
The interval is a diagnostic output, rather than an independently measured
guarantee of unique culprit identification.
Each stage has five seed fixtures with 100 repetitions per fixture.
Reported audit time covers the executed checks up to the first failure and
excludes model inference.

\noindent\textbf{Collusion controls.}
Table~\ref{tab:e-historical-collusion} reports three predefined controls,
with 500 repetitions each.
The $B+P$ control supplies conflicting provider-linked evidence while retaining
the original signed boundary records.
The $H+B$ control supplies an internally consistent seven-call transcript,
with no independent receipt recording the additional hidden execution.
The $H+B+P$ control supplies an internally consistent service-side transcript
without a conflicting external anchor.
The latter two controls are accepted and are not counted as blocked attacks.
These controls examine predefined evidence configurations rather than searching
all joint strategies available to colluding parties.
Witness expiry, late submission, and non-submission enforcement are outside
these measurements.

\begin{table*}[!t]
\centering
\caption{Historical collusion and observable-evidence controls. Acceptance in the last two rows is an evidence-boundary control, not a failed rejection of an observably inconsistent transcript.}
\label{tab:e-historical-collusion}
\footnotesize
\renewcommand{\arraystretch}{1.10}
\begin{tabular*}{\textwidth}{@{\hspace{4pt}}l@{\extracolsep{\fill}}p{.20\textwidth}p{.25\textwidth}p{.14\textwidth}p{.12\textwidth}@{\hspace{4pt}}}
\toprule
\textbf{Corrupt roles} & \textbf{Scenario} & \textbf{Independent evidence} & \textbf{Outcome} & \shortstack{\textbf{Responsibility}\\\textbf{interval}} \\
\midrule
$B+P$ & Cross-call provider-evidence splicing & Retained signed $H\!\rightarrow\!B$ boundary conflicts with provider-linked evidence & 500/500 rejected & $B$--$P$ \\
$H+B$ & One hidden execution omitted from an otherwise consistent transcript & No independent receipt for the hidden execution & 500/500 accepted & Unavailable \\
$H+B+P$ & Internally consistent service-side transcript & No conflicting external evidence anchor & 500/500 accepted & Unavailable \\
\bottomrule
\end{tabular*}
\end{table*}

